% =========================================================================
% theme/divs.tex
%
% LaTeX-окружения для семантических врезок учебника.
% Соответствуют fenced divs в Markdown-исходнике:
%
%   ::: opredelenie ... :::         → \begin{opredelenie} ... \end{opredelenie}
%   ::: napolyax ... :::             → \begin{napolyax} ... \end{napolyax}
%   ::: primer ... :::               → \begin{primer} ... \end{primer}
%   ::: kontrolnye-voprosy ... :::   → \begin{kontrolnyevoprosy} ... \end{kontrolnyevoprosy}
%
% Маппинг div-class → environment делается pandoc-filter'ом
% theme/filters/divs-to-env.lua (см. Спринт 1).
% =========================================================================

\usepackage[most]{tcolorbox}

% ---------- На полях. Вертикальная линейка слева, inline-заголовок ----------
\newtcolorbox{napolyax}[1][]{
  enhanced, blanker,
  borderline west={1.4pt}{0pt}{rulecol},
  left=12pt, right=4pt, top=2pt, bottom=2pt,
  fontupper=\small,
  before upper={\textbf{\small На полях.}\enspace},
  before skip=10pt,
  after skip=10pt,
  #1
}

% ---------- Определение. Двойная линейка сверху и снизу ----------
\newtcolorbox{opredelenie}[1][]{
  enhanced, blanker,
  borderline north={0.4pt}{0pt}{black},
  borderline north={0.4pt}{2pt}{black},
  borderline south={0.4pt}{0pt}{black},
  borderline south={0.4pt}{2pt}{black},
  left=4pt, right=4pt, top=6pt, bottom=6pt,
  fontupper=\small,
  before upper={\textbf{\small Определение.}\enspace},
  before skip=10pt,
  after skip=10pt,
  #1
}

% ---------- Пример. Тонкая линейка слева, меньший кегль ----------
\newtcolorbox{primer}[1][]{
  enhanced, blanker,
  borderline west={0.6pt}{0pt}{rulecol},
  left=10pt, right=4pt, top=2pt, bottom=2pt,
  fontupper=\footnotesize,
  before upper={\textbf{\footnotesize Пример.}\enspace},
  before skip=8pt,
  after skip=8pt,
  #1
}

% ---------- Контрольные вопросы — раздел в конце параграфа ----------
% Использование:
%   \begin{kontrolnyevoprosy}
%     \item Первый вопрос?
%     \item Второй?
%   \end{kontrolnyevoprosy}
\newenvironment{kontrolnyevoprosy}{%
  \par\vspace{14pt}
  \noindent{\bfseries Контрольные вопросы}\par
  \vspace{-2pt}
  \begin{enumerate}\setlength{\itemsep}{0pt}\setlength{\topsep}{2pt}%
}{%
  \end{enumerate}\par\vspace{6pt}%
}

% ---------- Задание — то же что контрольные, но другой заголовок ----------
\newenvironment{zadanie}{%
  \par\vspace{14pt}
  \noindent{\bfseries Задание}\par
  \vspace{-2pt}
  \begin{enumerate}\setlength{\itemsep}{0pt}\setlength{\topsep}{2pt}%
}{%
  \end{enumerate}\par\vspace{6pt}%
}

% ---------- Листинг 1С: рамка из двух горизонтальных линеек ----------
% Используется как обёртка для \input{listings/lst-N.tex} (output pygments)
\newtcolorbox{onecode}[1][]{
  enhanced,
  colback=white,
  colframe=white,
  frame hidden,
  borderline north={0.4pt}{0pt}{rulecol},
  borderline south={0.4pt}{0pt}{rulecol},
  arc=0pt,
  left=8pt, right=4pt, top=6pt, bottom=6pt,
  before skip=8pt,
  after skip=10pt,
  fontupper=\ttfamily\footnotesize,
  #1
}
